% Relation map for the controllability boundary: the bouncer at the one door (base)
% against the daemon at the mediation boundary (target). Arrows carry relations, not
% nouns. Hand-drawn in the pd figure language; no numeric data.
% [verified framework: Ramadge–Wonham controllability; event lists from the reference alphabet]
\begin{figure}[H]
\centering
\begin{tikzpicture}[pd figure,x=1cm,y=1cm]
  % panels
  \node[pd panel title,anchor=west] at (0,6.55) {Base: the club with one door};
  \node[pd panel title,anchor=west] at (8.3,6.55) {Target: the daemon at the mediation boundary};
  \draw[pd boundary] (0,0.2) rectangle (6.1,6.2);
  \draw[pd boundary] (8.3,0.2) rectangle (14.4,6.2);
  % base column
  \node[pd actor,text width=4.6cm] (bouncer) at (3.05,5.35) {the bouncer\\[-1pt]{\scriptsize refuses what passes the door}};
  \node[pd artifact,text width=4.6cm] (entry) at (3.05,3.55) {\textbf{entry}\\[-1pt]{\scriptsize mediated by the door: refusable}};
  \node[pd artifact,text width=4.6cm,fill=hhpaper] (thought) at (3.05,1.55) {\textbf{a patron's thought}\\[-1pt]{\scriptsize nothing the club owns mediates it}};
  \draw[pd arrow] (bouncer) -- (entry);
  \draw[pd guide] (bouncer) -- (thought);
  \node[pd note,anchor=west] at (0.15,2.6) {no arrow: cannot be refused};
  % target column
  \node[pd actor,text width=4.6cm] (daemon) at (11.35,5.35) {the daemon\\[-1pt]{\scriptsize refuses what crosses the boundary}};
  \node[pd artifact,text width=4.6cm] (sc) at (11.35,3.55) {$\Sigma_c$: \texttt{fs\_write}, \texttt{net\_egress}, \texttt{exec\_tool}, \texttt{git\_push}, \texttt{spawn\_child}, \texttt{api\_call}};
  \node[pd artifact,text width=4.6cm,fill=hhpaper] (su) at (11.35,1.55) {$\Sigma_u$: \texttt{model\_emit\_token}, \texttt{in\_context\_read}, \texttt{internal\_plan}, \texttt{hidden\_activation}};
  \draw[pd arrow] (daemon) -- (sc);
  \draw[pd guide] (daemon) -- (su);
  \node[pd note,anchor=west] at (8.45,2.6) {complete before the boundary sees them};
  % correspondence arrows, labelled by relation
  \draw[pd focus arrow,<->] (6.2,5.35) -- (8.2,5.35);
  \node[pd direct label,align=center] at (7.2,5.75) {stands at the only\\gate it owns};
  \draw[pd focus arrow,<->] (6.2,3.55) -- (8.2,3.55);
  \node[pd direct label,align=center] at (7.2,3.95) {refusable\\$\Leftrightarrow$ preventable};
  \draw[pd focus arrow,<->] (6.2,1.55) -- (8.2,1.55);
  \node[pd direct label,align=center] at (7.2,1.95) {unrefusable\\$\Leftrightarrow$ detect only};
  % criterion
  \node[pd note,anchor=north west,text width=14.2cm] at (0,0.05)
    {The map earns a prediction: a rule that mentions a thought but gates only the door stays preventable. In the criterion
     $\overline{K}\Sigma_u \cap \overline{L} \subseteq \overline{K}$ the door is $\Sigma_c$ and the thought is $\Sigma_u$.};
\end{tikzpicture}
\caption{What a runtime can prevent is decided by where an event crosses, not by what the rule is
about. Left, the base: a bouncer refuses entry and cannot refuse a thought. Right, the target: the
daemon refuses the six mediated effects and cannot refuse the four model-internal events. The arrows
between the columns carry the theorem: refusable at the boundary if and only if preventable; internal
if and only if detect-only. The event lists are the reference alphabet of the checker [internal,
\texttt{b3\_controllability.py}].}
\label{fig:swk-controllability-relation}
\end{figure}
